

Although these approaches have been useful in illustrating how identity affects elites' attitudes toward nuclear use \citep{tannenwald_nuclear_2007}, collateral damage, and other norm-breaking behavior, it has limited the scope of constructivist epistemology. The lack of quantitatively-driven, positivist study is striking given that experimental research has shown priming national identity can alter political behavior, such as by reducing affective polarization or changing attitudes toward immigration. Constructivist theorizing has often treated norms as being  internalized and difficult to contest. If it is the case that a norm can quickly be instantiated by identity-based primes, then that is proof of that norms' are weak rather than strong.  However, this outlook has served to limit the 